\scanpage{234}
\margin{- 189 - \qquad IX}
Donc les structures d'extension (14.1.2) et (14.1.1) peuvent
s'interpréter respectivement (VIII 3.7) comme équivalant à la donnée
d'homomorphisme
\[
(14.1.4)\qquad\underline M_K\xrightarrow{p_K}B_K,\quad
\underline M'_K\xrightarrow{p'_K}B'_K.
\]
Soient $W_K^{\natural}$ la biextension de $(B_K,B'_K)$ par $\mathbb G_{mK}$ qui définit la
dualité canonique entre ces schémas abéliens (8.2), et $E_K$ la biextension
de $(\underline M_K,\underline M'_K)$ par $\mathbb G_{mK}$ image inverse de $W_K^{\natural}$ par les homomorphismes (14.1.4).
Notons qu'on vérifie immédiatement que les trivialisations de $E_K$ correspondent
biunivoquement, de façon canonique, aux homomorphismes
\[
(14.1.5)\qquad i:\underline M_K\longrightarrow G_K
\]
qui relèvent $p_K:\underline M_K\longrightarrow B_K$, ou, symétriquement, aux homomorphismes
\[
(14.1.6)\qquad i':\underline M'_K\longrightarrow G'_K
\]
qui relèvent $P'_K:\underline M'_K\longrightarrow B'_K$.

Ceci posé, désignant par un exposant $^{\mathrm{an}}$ les groupes rigide-analytiques
sur $K$ associés aux groupes algébriques envisagés [13],
on définit dans [24] une suite exacte canonique
\[
(14.1.7)\qquad0\longrightarrow\underline M_K^{\mathrm{an}}
\xrightarrow{i^{\mathrm{an}}}G_K^{\mathrm{an}}\xrightarrow{e}A_K^{\mathrm{an}}\longrightarrow0,
\]
dans laquelle l'homomorphisme $i^{\mathrm{an}}$, qui peut aussi s'interpréter comme
un homomorphisme $i:\underline M_K\longrightarrow G_K$ pour des schémas en groupes ordinaires,
\uline{relève l'homomorphisme} $p$ (14.1.4), i.e. est de la forme envisagée dans
(14.1.5). Cette suite exacte est l'analogue rigide-analytique de la
suite exacte (13.3.4). Elle dépend de façon fonctorielle de $A_K$, et
\margin{507}
\margin{507}
\margin{233}
